\documentclass[11pt]{article}
\usepackage[margin=1in]{geometry}
\usepackage{amsmath,amssymb,amsthm,microtype,hyperref,booktabs}
\hypersetup{colorlinks=true,linkcolor=blue,citecolor=blue,urlcolor=blue}
\newtheorem{theorem}{Theorem}
\newtheorem{lemma}[theorem]{Lemma}
\newtheorem{corollary}[theorem]{Corollary}
\newcommand{\rev}{\operatorname{rev}}
\newcommand{\ds}{\operatorname{S}}
\newcommand{\zg}{\operatorname{Z}}
\title{Uniform digit-sum escape in reverse-and-add orbits}
\author{Working manuscript}
\date{October 14, 2026}
\begin{document}
\maketitle
\begin{abstract}
Fix an integer base $b\geq2$. For every positive integer, the base-$b$
digit sums along its indefinitely continued reverse-and-add orbit are
unbounded on every tail. An elementary carry argument controls new zero
runs in the symmetric digit-pair word: a long run either inherits an old
zero run or consumes a long block of coefficients equal to $b-1$.
Bounded digit sum would therefore imply bounded digit length. By extending
the dynamics to finite nonsymmetric profiles and inducting on their mass,
we prove an explicit escape bound depending only on the base and digit-sum
threshold, independent of the starting integer. We also give a square-root
lower bound for the running maximum of digit sums. In base ten its normalized lower limit is at least
$3/4$. The decimal unboundedness theorem, coarser explicit quantitative
bounds, and the stated uniform escape recurrence in base ten are
formalized in Lean 4.19 using its standard library.
\end{abstract}

\section{Statement, conventions, and prior work}
For a fixed base $b\geq2$, let
\[
 T_b(n)=n+\rev_b(n),\qquad n_t=T_b^t(n_0),\qquad n_0>0.
\]
Reversal uses the canonical base-$b$ expansion; leading zeros in the
reversed string have their usual numerical value. Write $\ds_b(n)$ for
the sum of its base-$b$ digits and $L(n)$ for its digit length. The orbit
continues after palindromes. For example, the hypothesis is not a rule
that stops the sequence when a palindrome is first reached.

\begin{theorem}\label{thm:main}
For every integer $b\geq2$, every $n_0>0$, and every $B,s\in\mathbb N$,
there exists a $t\geq s$ such that $\ds_b(n_t)>B$. Equivalently,
\[
 \limsup_{t\to\infty}\ds_b(n_t)=+\infty.
\]
\end{theorem}

The symmetric digit-pair representation is established prior art.
Styer's remainder-shift arrays encode a pair sum $q_i$ by its remainder
and quotient on division by the base \cite[Section 2]{styer}.
The half-pair-sum representation appears in Prosper and Veigneau's
pseudopalindromes \cite[Section 4.1]{pvpreprint}. Styer also records a
bound on residence at a fixed digit length \cite[p.~6]{styer}.
Explicit binary orbits with growing digit blocks were studied much
earlier by Trigg \cite{trigg}. Harborth \cite{harborth} constructed infinitely
many pairwise disjoint palindrome-free orbits in every power-of-two base;
their growing nonzero blocks have unbounded digit sum. None of these
ingredients is claimed here as a new invention.

The argument below isolates a dynamical bound on zero runs and deduces
unbounded digit sum and quantitative consequences. Our source comparison
uses Styer's manuscript and the 1997 Prosper--Veigneau preprint; the final
2001 publication \cite{pvpublished} was identified bibliographically but
its full text has not been inspected. Priority for the present conclusions
remains under review. This manuscript neither proves nor disproves the
existence of a decimal Lychrel number, and makes no claim that 196 avoids
palindromes. It does not assert the stronger pointwise limit
$\ds_b(n_t)\to\infty$.

\section{Exact carries and zero-run propagation}
Let $n=\sum_{i=0}^{L-1}d_i b^i$, where $0\leq d_i<b$,
$d_{L-1}>0$, and $L\geq1$. All indices increase from the units position.
Set
\[
 q_i=d_i+d_{L-1-i}\quad(0\leq i<L),\qquad
 M=\sum_{i=0}^{L-1}q_i=2\ds_b(n).
\]
The word $q$ is symmetric, $0\leq q_i\leq2b-2$, and its endpoints are
positive. Ordinary addition supplies carries $a_i\in\{0,1\}$ and digits
$r_i\in\{0,\ldots,b-1\}$ satisfying
\begin{equation}\label{eq:forward}
 a_0=0,\qquad q_i+a_i=r_i+b a_{i+1}.
\end{equation}
Put $h=a_L$, $c_i=a_{L-i}$, and $u_i=r_{L-1-i}$. Symmetry implies
\begin{equation}\label{eq:backward}
 q_i+c_{i+1}=u_i+b c_i.
\end{equation}
The output has canonical length $L+h$. In the nonoverflow case its
leading digit is positive, since $q_{L-1}>0$ and
$q_{L-1}+a_{L-1}<b$. Thus the next symmetric word is exactly
\begin{align}
 h=0:\quad &q'_i=r_i+u_i &&(0\leq i<L),\label{eq:next0}\\
 h=1:\quad &q'_0=q'_L=r_0+1,\qquad
 q'_i=r_i+u_{i-1} &&(1\leq i<L).\label{eq:next1}
\end{align}
Both formulas allow a zero units digit in the input. They also cover
$L=1$, with an empty interior range in \eqref{eq:next1}.

\begin{lemma}\label{lem:local}
If $r_i=u_i=0$ on an interval containing at least two consecutive
positions, then $q$ is identically $0$ or identically $b-1$ on that interval.
\end{lemma}
\begin{proof}
At each such position, \eqref{eq:forward}--\eqref{eq:backward} give
\[
 b(a_{i+1}-c_i)=a_i-c_{i+1}.
\]
The right side belongs to $\{-1,0,1\}$, so both sides vanish because
$b\geq2$. Hence $a_{i+1}=c_i$ and $a_i=c_{i+1}$. At adjacent positions
this implies $a_{i+2}=a_i$. If $a_i\ne a_{i+1}$, one of
$q_i=b a_{i+1}-a_i$ and $q_{i+1}=b a_{i+2}-a_{i+1}$ would be negative.
Therefore the three carries agree and
$q_i=q_{i+1}=(b-1)a_i$. Apply this to successive adjacent pairs.
An isolated simultaneous zero can instead have $q_i=b$; this explains
the restriction on interval length.
\end{proof}

Let $\zg(q)$ denote the greatest length of a consecutive zero run in $q$,
with value zero if $q$ has no zero entries.
\begin{lemma}\label{lem:gaps}
One reverse-and-add step satisfies
\begin{align}
 h=0:\quad \zg(q')&\leq
 \max\left\{\zg(q),1,\left\lfloor\frac{M}{b-1}\right\rfloor\right\},
 \label{eq:gap0}\\
 h=1:\quad \zg(q')&\leq
 \max\left\{\max(\zg(q)-1,0),2,
                   \left\lfloor\frac{M}{b-1}\right\rfloor+1\right\}.
 \label{eq:gap1}
\end{align}
\end{lemma}
\begin{proof}
Without overflow, a new zero run of length $k\geq2$ makes $r$ and $u$
zero on the same interval. Lemma~\ref{lem:local} says that this old
interval is all zero, giving $k\leq\zg(q)$, or all $b-1$, giving
$(b-1)k\leq M$. Runs of length one are already covered.

With overflow, consider a new zero run $[A,D]$ of length $k\geq3$.
Positive endpoints give $1\leq A\leq D\leq L-1$.
Equation~\eqref{eq:next1} gives $r=0$ on $[A,D]$ and $u=0$ on
$[A-1,D-1]$. Their common core $[A,D-1]$ has length $k-1\geq2$.
If its old coefficients are all $b-1$, then $(b-1)(k-1)\leq M$.

Otherwise the core is all zero. Its carry equations force
$a_A=a_D=c_A=c_D=0$. At the left adjacent column,
$u_{A-1}=c_A=0$ gives $q_{A-1}=b c_{A-1}$, while $a_A=0$ gives
$q_{A-1}+a_{A-1}<b$. Hence $q_{A-1}=0$.
At the right adjacent column, $r_D=a_D=0$ gives $q_D=b a_{D+1}$,
while $c_D=0$ gives $q_D+c_{D+1}<b$. Hence $q_D=0$ as well.
The old zero interval therefore extends to $[A-1,D]$, of length $k+1$.
Runs of length at most two complete the bound. In particular, long gaps
cannot arise by arbitrarily merging unrelated old zero runs.
\end{proof}

\section{Unbounded digit sums}
\begin{proof}[Proof of Theorem~\ref{thm:main}]
Suppose first that $\ds_b(n_t)\leq B$ for every $t\geq0$.
Since $n_t>0$, necessarily $B\geq1$. For the corresponding words $q_t$,
let $L_t=L(n_t)$ and set
\begin{equation}\label{eq:KH}
 K_B=\max\left\{\zg(q_0),2,
                  \left\lfloor\frac{2B}{b-1}\right\rfloor+1\right\},
 \qquad H_B=2B+(2B-1)K_B.
\end{equation}
The mass is at most $2B$, so Lemma~\ref{lem:gaps} gives
$\zg(q_t)\leq K_B$ inductively. If $m_t$ is the number of nonzero
coefficients, then $1\leq m_t\leq2B$. Positive endpoints leave at most
$m_t-1$ intervening zero runs. Consequently
\begin{equation}\label{eq:length}
 L_t\leq m_t+(m_t-1)K_B\leq H_B.
\end{equation}
But $\rev_b(n_t)\geq1$, so $n_t\geq n_0+t$. This contradicts
$n_t<b^{H_B}$ for all $t$. Starting the same argument at $n_s$ proves
the assertion on every tail. The case $B=0$ is immediate.
\end{proof}

The proof yields unbounded nonzero-digit count as well: if $w_b(n)$
counts nonzero base-$b$ digits, then
$\ds_b(n)\leq(b-1)w_b(n)$. It also excludes every positive orbit with
eventually constant or eventually bounded digit sum.

\section{Explicit quantitative bounds}
\begin{lemma}\label{lem:width}
Let $r_b$ be the least integer $r$ with $b\leq2^r$, and put
$\kappa_b=r_b+1$. Every positive orbit satisfies
\[
 L_t\geq L_0+\left\lfloor\frac{t}{\kappa_b}\right\rfloor.
\]
\end{lemma}
\begin{proof}
Write $x=q_0>0$ for the outer coefficient and $e=a_{L-1}\in\{0,1\}$.
Without overflow, $x+e<b$ and \eqref{eq:next0} gives
$x'=2x+e\geq2x$. If $r_b+1$ consecutive steps had no overflow,
the first $r_b$ would make the outer coefficient at least $2^{r_b}\geq b$,
forcing overflow at the next step. Thus each disjoint block of
$\kappa_b$ steps has an overflow. Length increases by exactly one at
an overflow and never decreases.
\end{proof}

\begin{corollary}\label{cor:escape}
If $\ds_b(n_0)\leq B$, define $K_B,H_B$ by \eqref{eq:KH}.
Then some time satisfies
\[
 0\leq t\leq\kappa_b(H_B-L_0+1),\qquad \ds_b(n_t)>B.
\]
For fixed $b,n_0$, the running maximum
$D_N=\max_{0\leq t\leq N}\ds_b(n_t)$ satisfies
$D_N=\Omega(\sqrt N)$.
\end{corollary}
\begin{proof}
The assumption implies $B\geq1$ and $H_B\geq L_0$ by the counting
argument at time zero. If the bound on digit sums held through
$E=\kappa_b(H_B-L_0+1)$, the same finite induction would give
$L_E\leq H_B$. Lemma~\ref{lem:width} instead gives $L_E\geq H_B+1$.
If the initial digit sum already exceeds $B$, take $t=0$.

For the second assertion, use $B=D_N$ on the finite prefix to obtain
\[
 L_0+\left\lfloor\frac N{\kappa_b}\right\rfloor
 \leq2D_N+(2D_N-1)
 \max\left\{\zg(q_0),2,
             \left\lfloor\frac{2D_N}{b-1}\right\rfloor+1\right\}.
\]
The right side is $O(D_N^2)$, with constants depending only on the
fixed base and initial word. This proves the assertion.
\end{proof}

\subsection*{A sharper decimal constant}
For $b=10$ an outer-column potential improves the width estimate.
Define $p$ on $\{1,\ldots,18\}$ by
\[
\begin{array}{c|ccccc}
x&1&2&3,4,5&6,7,8,9,10&11,\ldots,18\\ \hline
p(x)&4&3&2&1&0.
\end{array}
\]
Let $x=q_0$, $e=a_{L-1}$, and $h=a_L$. If $h=0$, then
$x'=2x+e$ with $x+e<10$. If $h=1$, then
$x'=x+1-10\lfloor x/10\rfloor$ with $x+e\geq10$.
Checking these two formulas against the five intervals of the table gives
\[
 p(x')+1\leq p(x)+4h.
\]
For clarity, when $h=0$, $x=1$ leads to $x'=2$ or $3$;
$x=2$ to $4$ or $5$; $3\leq x\leq5$ to $6\leq x'\leq11$;
and $6\leq x\leq9$ to $12\leq x'\leq18$.
When $h=1$, the cases $x=9$, $x=10$, and $11\leq x\leq18$
give $x'=10$, $x'=1$, and $2\leq x'\leq9$, respectively.
These ranges prove the inequality directly. Summing along the orbit yields
\begin{equation}\label{eq:decimalwidth}
 t+4L_0\leq4L_t+p(q_0)\leq4L_t+4.
\end{equation}
This argument is also checked in the accompanying Lean source.

Thus the decimal escape horizon in Corollary~\ref{cor:escape} improves to
$4(H_B-L_0)+5$. More explicitly, put $A=\max\{\zg(q_0),2\}$.
For $D=D_N$, the sharp gap bound obeys
$K_D\leq A+\lfloor2D/9\rfloor$, and therefore
\[
 H_D\leq\frac49D^2+2(A+1)D.
\]
Combining with \eqref{eq:decimalwidth} and completing the square gives
\begin{equation}\label{eq:explicit}
 D_N\geq\frac34\left[
 \sqrt{N+4L_0-4+9(A+1)^2}-3(A+1)\right].
\end{equation}
The radicand is positive because $L_0\geq1$. In particular,
\begin{equation}\label{eq:liminf}
 \liminf_{N\to\infty}\frac{D_N}{\sqrt N}\geq\frac34.
\end{equation}
Equations~\eqref{eq:explicit}--\eqref{eq:liminf} concern the running
maximum, not each individual digit sum. For example,
$1099+9901=11000$ decreases the digit sum from $19$ to $2$.

\section{Uniform escape independent of the starting integer}
The preceding quantitative estimates depend on the initial word. The next
result removes that dependence, at the cost of a much larger bound.
For fixed $b$, let $\kappa_b$ be as in Lemma~\ref{lem:width}, and define
\begin{equation}\label{eq:recursion}
 E_b(1)=1,\qquad
 E_b(M)=\kappa_b\bigl[2(M-1)E_b(M-1)+2M-1\bigr]\quad(M\geq2).
\end{equation}
\begin{theorem}\label{thm:uniform}
For every $b\geq2$, $B\geq1$, $n_0>0$, and $s\geq0$, there exists
$0\leq j\leq E_b(2B)$ such that
\[
 \ds_b(T_b^{s+j}(n_0))>B.
\]
Thus the waiting-time bound depends only on $b$ and $B$, not on the
starting integer, its digit length, or the starting time.
\end{theorem}
The bound is deliberately coarse. For example,
$E_{10}(2)=25$, $E_{10}(3)=525$, and $E_{10}(4)=15785$.
These are upper bounds, not optimal escape times.

\subsection*{Finite profiles without symmetry}
Consider any finitely supported integer sequence $q=(q_i)_{i\geq0}$ with
$0\leq q_i\leq2b-2$ and $q_0>0$. Write
$M(q)=\sum_iq_i$, $R(q)=\max\{i:q_i>0\}$, and let $Z(q)$ be the
greatest length of a zero run within $[0,R(q)]$, with value zero if none
exists. The infinite zero tail is excluded.
Define forward carries and digits by \eqref{eq:forward}, now for all
$i\geq0$. Independently define backward carries and digits by
\eqref{eq:backward}, setting $c_i=0$ sufficiently far to the right and
scanning toward zero. Put $h=c_0$ and define
\begin{equation}\label{eq:halfline}
 \begin{aligned}
 h=0:&\quad \Phi(q)_i=r_i+u_i &&(i\geq0),\\
 h=1:&\quad \Phi(q)_0=r_0+1,\qquad
                 \Phi(q)_i=r_i+u_{i-1} &&(i\geq1).
 \end{aligned}
\end{equation}
The forward terminal digit $r_{R(q)+1}$ is included. The image again has
finite support, coefficient bound $2b-2$, and positive zeroth entry.
In the $h=1$ branch the boundary coefficient is at most $b\leq2b-2$.
In the other branch it is $2q_0+c_1>0$.
For a genuine symmetric pair-sum word, extended by zeros, $\Phi$ is
exactly the ordinary next-pair-sum operation. No limiting or infinite
integer is used in this extension.

Lemma~\ref{lem:local} used only the two carry equations, so its proof
remains valid without symmetry. The same holds for the gap inequalities
of Lemma~\ref{lem:gaps}. To check their right boundary, observe first that
support never retreats. If $h=0$, either $a_{R+1}=1$, producing a
nonzero digit at $R+1$, or $r_R=q_R+a_R>0$. If $h=1$, then
\[
 \Phi(q)_{R+1}=a_{R+1}+u_R>0:
\]
when the first term vanishes, $1\leq q_R<b$ and $u_R=q_R>0$.
Thus $R(\Phi(q))=R(q)+1$ when $h=1$, and
$R(\Phi(q))\geq R(q)$ otherwise. No new zero interval can extend into
the infinite tail, because that tail starts after the positive last
coefficient. The long-core and boundary-extension arguments consequently
apply exactly as before.

At a step with $h=0$, the zeroth coefficient at least doubles. Hence
$\kappa_b$ such steps cannot occur consecutively, and for an orbit of
profiles
\begin{equation}\label{eq:profilegrowth}
 R(q_t)\geq R(q_0)+\left\lfloor\frac t{\kappa_b}\right\rfloor.
\end{equation}
If its mass is at most $M$ and its internal gaps are at most $K$, the
usual count gives
\begin{equation}\label{eq:profilewidth}
 R(q_t)+1\leq M+(M-1)K.
\end{equation}

\subsection*{Isolation by a long zero gap}
Suppose $q=p+v$, where $p_0>0$, both parts are nonzero, and
\[
 A=\max\operatorname{supp}(p),\quad
 D=\min\operatorname{supp}(v),\quad g=D-A-1\geq2.
\]
A zero coefficient resets a carry passing through it. Forward normalization
of $p$ can extend its support only to $A+1$, and backward normalization
of $v$ only to $D-1$. The gap therefore prevents either directional
carry from reaching the other component. Both digit normalizations of
$p+v$ are the sums of the corresponding separate normalizations, and its
boundary carry $h$ equals that of $p$.
Define, with $u_i(v)=0$ for $i<0$,
\[
 \Psi_h(v)_i=r_i(v)+u_{i-h}(v)\qquad(i\geq0).
\]
There is no extra boundary digit in this right component. Equation
\eqref{eq:halfline} gives the exact identity
\begin{equation}\label{eq:split}
 \Phi(p+v)=\Phi(p)+\Psi_h(v).
\end{equation}
The support of $\Phi(p)$ ends no later than $A+1$, while that of
$\Psi_h(v)$ begins no earlier than $D-1$. Thus the gap shrinks by at
most two positions per step. The right component stays nonzero because
its forward digits represent the positive integer $\sum_i v_i b^i$,
and $\Psi_h(v)$ is coefficientwise at least that nonzero digit sequence.
Consequently, if $g\geq2T+2$, identity \eqref{eq:split} can be iterated
through time $T$: the components stay disjoint, the left one evolves
autonomously by $\Phi$, and the right one always contributes at least
one unit of mass.

\begin{proof}[Proof of Theorem~\ref{thm:uniform}]
We prove by induction on $M\geq1$ that no profile can have mass at most
$M$ at all times $0,\ldots,E_b(M)$. For $M=1$, the only possible initial
profile is $q_0=1$, $q_i=0$ for $i>0$; its next mass is two.

For $M\geq2$, set $T=E_b(M-1)$, $K=2T+1$, and
\[
 H=M+(M-1)K=2(M-1)T+2M-1.
\]
Suppose that all masses through time $E_b(M)=\kappa_b H$ are at most
$M$. This horizon is at least $T$. If the initial profile has a zero gap
of length at least $2T+2$, split at that gap. The isolation argument gives
an autonomous left profile whose mass is at most $M-1$ through time $T$,
since the right component is nonzero throughout. This contradicts the
induction hypothesis.

Otherwise the initial gaps are at most $K$. Recurrence
\eqref{eq:recursion} gives $T\geq M-1$, and hence
\[
 K=2T+1\geq2M-1\geq M+1
       \geq\left\lfloor\frac M{b-1}\right\rfloor+1,
 \qquad K\geq2.
\]
The gap inequalities preserve the bound $K$ throughout the assumed
prefix. Equation~\eqref{eq:profilewidth} bounds its support width by
$H$, while \eqref{eq:profilegrowth} at time $\kappa_b H$ makes the
right endpoint at least $H$, a contradiction. This proves the induction.

A genuine pair-sum profile has mass $2\ds_b(n)$. Apply the result with
$M=2B$ to the profile of $T_b^s(n_0)$. If its initial mass already exceeds
$2B$, take $j=0$; otherwise the induction supplies the required later time.
\end{proof}

For completeness, induction on \eqref{eq:recursion} also gives the
simpler closed upper bound
\[
 E_b(M)\leq(4\kappa_b)^{M-1}M!.
\]
Indeed, for $M\geq2$ put $U=(4\kappa_b)^{M-2}(M-1)!\geq1$.
The available margin after the induction term is
$[4\kappa_b M-2\kappa_b(M-1)]U=2\kappa_b(M+1)U$,
which exceeds $\kappa_b(2M-1)$.

\begin{corollary}
Under the hypotheses of Theorem~\ref{thm:uniform}, put $F=E_b(2B)$ and
$A_B=\{t\geq0:\ds_b(T_b^t(n_0))>B\}$. Every interval of $F+1$
consecutive integer times meets $A_B$. In particular,
\[
 \liminf_{N\to\infty}\ \inf_{s\geq0}
 \frac{\#(A_B\cap\{s,\ldots,s+N-1\})}{N}
 \geq\frac1{F+1}.
\]
\end{corollary}
\begin{proof}
Apply Theorem~\ref{thm:uniform} at each starting time and partition any
interval of length $N$ into $\lfloor N/(F+1)\rfloor$ disjoint complete
blocks of length $F+1$.
\end{proof}
This bounded-waiting-time conclusion still permits isolated small digit
sums at arbitrarily late times. It therefore does not establish pointwise
divergence of digit sums.

\section{Formal verification and scope}
The standalone file \texttt{LychrelDigitSums.lean} uses Lean 4.19.0 and
\texttt{Std}, without mathlib. It includes the canonical decimal digit
representation, the exact connection from the carry-word dynamics to
ordinary integer iteration, zero-run propagation, and the unboundedness
theorem on every tail. The principal declaration is
\begin{center}
\texttt{DecimalUnboundedMass.digitSum\_unbounded\_after}.
\end{center}
No unproved local propagation premise remains in that theorem.

The quantitative extension defines a computable decimal length and an
attained finite running maximum. For every $n>0$ and $B\in\mathbb N$, put
\[
 L=L(n),\qquad \overline K=\max\{L,2B+2\},\qquad
 \overline H=(2B+1)(\overline K+1).
\]
The checked theorem \texttt{DecimalQuantitativeMass.digitSum\_escape}
proves
\[
 \exists\,t\leq4(\overline H-L)+5:\quad \ds_{10}(T_{10}^t(n))>B.
\]
This statement needs no initial digit-sum assumption. The file also
formalizes a running-maximum inequality. These quantitative estimates use
coarser constants than the written proof.

\subsection*{Completed uniform decimal formalization}
Nine additional modules in the accompanying \texttt{uniform\_lean}
directory formalize the nonsymmetric profile map, its exact support,
zero-gap propagation, finite-time isolation, induction on mass, and the
embedding of ordinary integer orbits. No isolation or dynamical closure
law is assumed. The checked recurrence is exactly \eqref{eq:recursion}
for $b=10$, with an additional harmless definition at zero:
\[
 E(0)=0,\quad E(1)=1,\qquad
 E(M)=5\bigl[2(M-1)E(M-1)+2M-1\bigr]\quad(M\geq2).
\]
\begin{samepage}
The declaration
\begin{center}
\texttt{DecimalUniformEscape.digitSum\_uniform\_window\_sharp}
\end{center}
proves, for every $n>0$ and $B,s\in\mathbb N$,
\[
 \exists t\in\{s,\ldots,s+E(2B)\}:
       \ds_{10}(T_{10}^t(n))>B.
\]
\end{samepage}
Thus the explicit decimal bound in Theorem~\ref{thm:uniform} is formalized.
Its deadline depends on $B$ alone, not on the initial number, its decimal
length, or the starting time $s$. The statement includes $B=0$ and inputs
already exceeding the threshold. A separate checked specialization applies
it to $n=196$ without asserting palindrome avoidance.

The supporting profile theorem is
\texttt{DecimalProfile.sharp\_uniform\_escape}. Its proof uses the direct
mass-one base case, the positive-endpoint width estimate, and the sharper
zero-gap propagation bound. The new map retains the terminal forward carry
even for nonsymmetric profiles. The file \texttt{DecimalEmbedding.lean}
proves exact correspondence with ordinary reverse-and-add at every step
and identifies profile mass with twice the ordinary digit sum. These are
concrete theorems, not hypotheses of the final result. The initial proof
with a coarser uniform recurrence is also retained in the source package.

All nine new modules and the existing standalone dependency compile in
Lean 4.19.0. The final declarations use only \texttt{propext},
\texttt{Classical.choice}, and \texttt{Quot.sound}; they use no custom
axioms, admitted proofs, or native computation oracle. The general-base
extension and the sharper decimal constant $3/4$ remain written,
independently audited conclusions rather than claims about these Lean
declarations. Source, build instructions, and axiom-audit records accompany
this draft.

The result excludes the previously investigated fixed-distance wall
construction, whose all-time finite embedding would force constant
positive digit sum. A successful decimal Lychrel construction must allow
unbounded digit sum. Unbounded digit sum itself does not imply palindrome
avoidance: the theorem also applies to orbits that have reached palindromes.

\paragraph{Preparation and status.}
OpenAI language models were used substantially in developing arguments,
writing and checking Lean code, preparing this text, and assisting source
searches. The mathematical audits in the project are AI-assisted internal
checks, not external human peer review. No journal submission, acceptance,
or claim of publication priority is represented by this draft.

\begin{thebibliography}{9}
\bibitem{styer}
R. Styer, \emph{The Palindromic Conjecture and The Fibonacci Sequence},
Villanova University manuscript, revised September 1986, 11 pages.
\url{https://www41.homepage.villanova.edu/robert.styer/PalindromePaper1986.pdf}.

\bibitem{pvpreprint}
V. Prosper and S. Veigneau, \emph{On the palindromic reversal process},
preprint dated June 18, 1997, 13 pages.
\url{https://www.p196.org/files/revpal.pdf}.

\bibitem{pvpublished}
V. Prosper and S. Veigneau, On the palindromic reversal process,
\emph{Calcolo} \textbf{38} (2001), 129--140.
\url{https://doi.org/10.1007/s100920170002}.
Bibliographic record checked; final full text not inspected.

\bibitem{trigg}
C. W. Trigg, Versum sequences in the binary system,
\emph{Pacific Journal of Mathematics} \textbf{47} (1973), no.~1, 263--275.
\url{https://msp.org/pjm/1973/47-1/pjm-v47-n1-p25-p.pdf}.

\bibitem{harborth}
H. Harborth, On palindromes,
\emph{Mathematics Magazine} \textbf{46} (1973), no.~2, 96--99.
\url{https://doi.org/10.1080/0025570X.1973.11976286}.
\end{thebibliography}
\end{document}
